\documentclass{beamer}
% Für schnelleres Kompilieren:
%\documentclass[draft]{beamer}
%\includeonlyframes{toc,current}

\mode<presentation>
{
  \usetheme{Berlin}
  % Pittsburgh, Dresden, Boadilla, Berlin, Goettingen, Luebeck, Rochester
  \setbeamercovered{transparent}
}

\usepackage[german]{babel}
\usepackage[utf8]{inputenc}
\usepackage[T1]{fontenc}
\usepackage{lmodern}
%\usepackage{times}
\usepackage{multimedia}

\title[GRAL WS 06/07]{Nichtrealistische Computergraphik}
\author[Kungel, Nägele, Tammen]{Christian Kungel, Kornelius Nägele, Jan Tammen}

\institute{
  Seminar Graphische Algorithmen WS 06/07\\
  Fakultät Informatik,\\
  HTWG Konstanz
}

%%% Informationen ueber das PDF-Dokument setzen
\hypersetup{
  pdftitle={Nichtrealistische Computergraphik},%
  pdfauthor={Jan Tammen,%
  			 Kornelius Nägele,%
  			 Christian Kungel},%
  pdfsubject={Graphische Algorithmen},%
  pdfcreator={LaTeX with pdfeTeX},%
  pdfproducer={LaTeX with pdfeTeX},%
  pdfkeywords={Nichtrealistische Computergraphik, NPR},
  % pdfpagelayout=TwoColumnRight,
  pdffitwindow=true,
  pdfpagemode=FullScreen,
%  pdfstartview=FitH,
}

\date{\today}
\subject{Nichtrealistische Computergraphik}

\pgfdeclareimage[height=1cm]{htwg-logo}{../images/htwg-logo}
\logo{\pgfuseimage{htwg-logo}}

% Folgendes sollte gelöscht werden, wenn man nicht am Anfang jedes
% Unterabschnitts die Gliederung nochmal sehen möchte.
\AtBeginSubsection[]
{
 \begin{frame}<beamer>
   \frametitle{Gliederung}
   \tableofcontents[currentsection,currentsubsection]
 \end{frame}
}

% Falls Aufzählungen immer schrittweise gezeigt werden sollen, kann
% folgendes Kommando benutzt werden:
% mit highlight: <+- | alert@+>
\beamerdefaultoverlayspecification{<+->}

\begin{document}

% Frame: Titelseite
\begin{frame}
  \titlepage
\end{frame}

% Frame: Gliederung
\begin{frame}[label=toc]
  \tableofcontents
\end{frame}

\section{Einleitung}
\begin{frame}{Mozart - Einleitung}
  \begin{itemize}
    \item Mozart implementiert die multiparadigmische Programmiersprache Oz
    \item Forschungsprojekt seit ca.\ 1995 (B, D, S)
    \item Oz ist plattformunabhängig (Byte-Code ähnlich wie Java)
    \item Open-Source-Lizenz
  \end{itemize}
\end{frame}

\begin{frame}{Mozart - Features}
  \begin{itemize}
    \item Unterstützung mehrerer Programmierparadigmen
    \item Nebenläufigkeit (Threads, Synchronisation)
    \item Transparente Verteilung
    \item Dynamische Typisierung
  \end{itemize}
\end{frame}

\begin{frame}{Oz - Programmiermodell (OPM)}
  \begin{itemize}
    \item Berechnungen laufen in sog.\ "`Spaces"' (Berechnungsraum) ab
    \item \textsl{Aktoren} kommunizieren darin über einen gemeinsamen 
    \textsl{Speicher}, können Information ablegen und abrufen
    \item Grundlage von Oz: Concurrent Constraint Programming, Erweiterungen 
    durch \textsl{syntactic sugar}
  \end{itemize}
\end{frame}

\begin{frame}{Oz - Datentypen}
  \begin{itemize}
    \item Basisdatentypen: \texttt{Number: Int, Float, Char}
    \item Zusammengesetzte Typen: \texttt{Record: Tuple, Literal, Bool}
    \item \texttt{Chunk}: erlaubt Erstellung abstrakter Datentypen, vordefiniert
    z.\,B.\ \texttt{Array, Dictionary, Class}
    \item \texttt{Thread, Space, ByteString, \ldots}
  \end{itemize}
\end{frame}

\section{Grundlegende Techniken}
\begin{frame}{Einsatzebenen}
  \textbf{Non-Photorealistic Rendering (NPR)} kann auf drei Ebenen eingesetzt werden:
  \begin{itemize}
    \item Geometrieraum
    \item Projektionsraum
    \item Bildraum
  \end{itemize}
\end{frame}

\subsection{Techniken im Geometrieraum}
\begin{frame}{Techniken im Geometrieraum}
  \begin{itemize}
    \item direkter Einfluss auf die 3D-Geometrie der Objekte
    \item Nachteil: erfordert einen hohen Zeitaufwand
    \item wird daher selten bevorzugt
    \item Techniken:
    \begin{itemize}
      \item Computer Sketching
      \item Noise Modifier
      \item Level of Detail
    \end{itemize}
  \end{itemize}
\end{frame}

\begin{frame}{Computer Sketching}
  \begin{itemize}
    \item Paul Bourk, 1990
    \item Frage: Können CAD-Anwendungen auch unscharfe,
    an menschliche Skizzen erinnernde Entwürfe
    ausgeben?
    \item Kriterien für Skizzen:
    \begin{itemize}
      \item Linien enden nicht exakt an Schnittpunkten (Kreuzungen)
      \item Linien sind nicht gradlinig sondern "`verwackelt"'
    \end{itemize}
  \end{itemize}
\end{frame}

\begin{frame}{Computer Sketching}
  \begin{itemize}
    \item Hauptansatz: direkte Manipulation der 3D-Geometrie
    \item Umsetzung:
    \begin{itemize}
      \item Linien werden um randomisierte Faktoren verlängert
      \item Aufspalten der Linien und Mittelpunktverschiebung
    \end{itemize}
  \end{itemize}
\end{frame}

\begin{frame}{Computer Sketching}
  \includegraphics[width=0.5\textwidth]{../images/haus.jpg}
\end{frame}

\begin{frame}{Noise Modifier}
  \begin{itemize}
    \item "`Rauschen"' in 3D-Modellen
    \item Geometriepunkte werden um einen zufälligen Faktor
    verschoben
    \item meist Anwendung auf glatte Flächen, um
    "`natürliche"' Flächen zu erzeugen (Landschaften)
  \end{itemize}
\end{frame}

\begin{frame}{Level of Detail}
  \begin{itemize}
    \item Einsatz in vielen Bereichen:
    \begin{itemize}
      \item Kartographie: angezeigte Detailstufe in Abhängigkeit vom
      Maßstab (Generalisierung)
      \item photorealistische Darstellungen: angezeigte Detailstufe der
      Szene in Abhängigkeit von der Betrachterposition
      \item NPR: Detailstufe in Abhängigkeit zur "`Wichtigkeit"' der
      Objekte
      \begin{itemize}
        \item Eine Möglichkeit ist die Meshsimplifizierung
      \end{itemize}
    \end{itemize}
  \end{itemize}
\end{frame}

\begin{frame}{Level of Detail}
  \begin{itemize}
    \item Meshsimplifizierung:
    \begin{itemize}
      \item reduziert die Anzahl der Polygone eines 3D Modells
      \item äußere Form wird dabei approximiert
      \item Dadurch: schnellere Verarbeitung in Echtzeitsystemen möglich
    \end{itemize}
  \end{itemize}
\end{frame}

\subsection{Techniken im Projektionsraum}
\begin{frame}{Techniken im Projektionsraum}
  \begin{itemize}
    \item meistens Kombination aus Projektions- und
    Bildraumverfahren
    \item dadurch leichteres Arbeiten mit 3D-Modellen und
    Erhöhung der grafischen Qualität der Ausgabe
    \item Techniken:
    \begin{itemize}
      \item G-Buffer
      \item Line-Rendering
      \item Stroke Textures
    \end{itemize}
  \end{itemize}
\end{frame}

\begin{frame}{G-Buffer}
  \begin{itemize}
    \item "`Geometrie-Buffer"' Saito \& Takahashi, 1990
    \item in den meisten kommerziellen Cartoon-Render-Paketen implementiert
    \item Prinzip: zusätzliche Speicherung von geometrischen
    Daten zur Informationsgewinnung im Bildraum
    \begin{itemize}
      \item Flächen-Normalen des 3D-Modells
      \item UV-Koordinaten des 3D-Modells
      \item Z-Buffer des 3D-Modells
    \end{itemize}
    \item Grundidee: Informationen zur Bilderzeugung aus dem G-Buffer
    extrahieren, analysieren und zusammenfügen.  
  \end{itemize}
\end{frame}

\begin{frame}{Line-Rendering}
  \begin{itemize}
    \item Simulation der Semantik von Linien
    \item geometrische Semantik:
    \begin{itemize}
      \item Begrenzungen, Silhouetten, Diskontinuitäten (Falten)
    \end{itemize}
    \item stilistische Semantik:
    \begin{itemize}
      \item Stärke der Linien, Transparenz, Linientyp (durchgezogen,
      gestrichelt, etc.)
    \end{itemize}
  \end{itemize}
\end{frame}

\begin{frame}{Line-Rendering}
  \begin{itemize}
    \item Abspeicherung dieser zusätzlichen Attribute in
    geeigneter Form
    \begin{itemize}
      \item Eine Möglichkeit: Speicherung in Matrixform
    \end{itemize}
    \item Generierung der Attribute
    \begin{itemize}
      \item z.B. durch Inferenzregeln zu Bestimmung der
      Charakteristika der einzelnen Linien
    \end{itemize}
    \item Gute Kombinationsmöglichkeiten mit anderen
    NPR-Techniken (etwa G-Buffer)
  \end{itemize}
\end{frame}

\begin{frame}{Stroke Textures}
  \begin{itemize}
    \item Winkenbach \& Salesin, 1994
    \item Prinzip basiert auf konventionellen 2D-Malsystem
    \begin{itemize}
      \item "`strichbasierte"' Texturen werden angefertigt
      \item mittels Texturemapping auf 3D-Objekte aufgebracht
      \item Dabei: Anwendung des LOD-Prinzips auf die Textur
    \end{itemize}
  \end{itemize}
\end{frame}

\begin{frame}{Stroke Textures}
  \includegraphics[width=\textwidth]{../images/strokes.jpg}
  \begin{itemize}
    \item Linienstärken und "`Detail Tags"' sind durch den
    Benutzer spezifizierbar
  \end{itemize}
\end{frame}

\subsection{Techniken im Bildraum}
\begin{frame}{Techniken im Bildraum}
  \begin{itemize}
    \item vorhandenes Bild wird analysiert und neu
    gezeichnet
    \item Techniken:
    \begin{itemize}
      \item Texture Elements
      \item Hairy brushes
      \item Iterated Function System (IFS)
    \end{itemize}
  \end{itemize}
\end{frame}

\begin{frame}{Texture Elements}
  \begin{itemize}
    \item Mezei et al., 1974
    \item Ursprünglich: Synthese sehr "`naturalistischer"' Bilder
    \item Beobachtung: in der Natur kommen sich oft
    wiederholende Muster vor
  \end{itemize}
\end{frame}

\begin{frame}{Texture Elements}
  \begin{itemize}
    \item Umsetzung: vorgefertigte grafische Subelemente
    werden zusammengefügt, um neue Bilder zu
    erzeugen
    \item Variationen durch Skalierung, Rotation und andere
    Verzerrungen
    \item Erzielt werden Verallgemeinerungen der
    Oberflächeneigenschaften, die wie menschliche
    Illustrationen wirken
  \end{itemize}
\end{frame}

\begin{frame}{Hairy Brushes}
  \begin{itemize}
    \item Strassmann, 1986
    \item Verhalten von Pinsel und Farbe soll simuliert werden
    \item Ansatz: Pinsel als 1D-Abdruck der Pinselborsten
    \item Borsten laufen an einer parametrischen Kurve
    entlang, die durch Knoten definiert ist
  \end{itemize}
\end{frame}

\begin{frame}{Hairy Brushes}
  \begin{itemize}
    \item Knoten enthalten Informationen über Position und
    Druck des Pinsels
    \item jeder "`Strich"' wird mit einer speziellen Menge an
    Farbe gezeichnet
    \item Menge nimmt mit der Länge der Linie ab
    \item erneutes "`Eintunken"' des Pinsels ist möglich
  \end{itemize}
\end{frame}

\begin{frame}{Iterated Function Systems (IFS)}
\begin{itemize}
 \item IFS: Systeme von Funktionen die wiederholt
 ausgeführt werden
 \item Einsatz: zur Beschreibung von Fraktalen
 \item Fraktale sind selbstbezüglich (rekursiv)
 \begin{itemize}
   \item Beispiel: Barnsley-Farn
     \begin{columns}
    \column{.2\textwidth}
	\includegraphics[width=\textwidth]{../images/BarnsleyFarn.jpg}
	\column{.2\textwidth}
    \includegraphics[width=\textwidth]{../images/BarnsleyFarnDetail.jpg}
  \end{columns}
 \end{itemize}
\end{itemize}

\end{frame}

\begin{frame}{Iterated Function Systems (IFS)}
  \begin{itemize}
    \item Nutzung des Prinzips zur fraktalen Kompression von
    Bildern (Barnsley, 1988)
    \item Kompressionstechnik basiert auf dem "`Collage-
    Theorem"' (Selbstähnlichkeit von Bildern)
  \end{itemize}
\end{frame}

\section{Einsatzgebiete}
\subsection{Einsatzgebiete}

\begin{frame}{Architektur}
  \begin{columns}
    \column{.6\textwidth}
	\includegraphics[width=0.9\textwidth]{../images/dach.jpg}
	\column{.4\textwidth}
	Abstrakte Landschaftsdarstellung in Bauplänen\\
	Quelle: Uni Rostock
  \end{columns}
\end{frame}

\begin{frame}{Medizin, Biologie}
  \begin{columns}
    \column{.6\textwidth}
	\includegraphics[width=\textwidth]<1>{../images/herz1.jpg}
	\includegraphics[width=\textwidth]<2->{../images/herz2.jpg}
	\column{.4\textwidth}
	Abstrakte Darstellung von Organismen\\
	Quelle: K. Hartmann, Magdeburg
  \end{columns}
\end{frame}

\begin{frame}{Spiele}
  \begin{columns}
    \column{.6\textwidth}
	\includegraphics[width=\textwidth]{../images/Zelda.jpg}
	\column{.4\textwidth}
	Echtzeitdarstellung von comichaften Umgebungen\\
	Zelda, ein Nintendo-Spiel\\
	Quelle: Wikipedia
  \end{columns}
\end{frame}

\begin{frame}{Kunst}
  \begin{columns}
    \column{.6\textwidth}
	\includegraphics[width=\textwidth]{../images/kunscht.jpg}
	\column{.4\textwidth}
	Hervorhebung durch Verfälschung von Fotos\\
	Quelle: Uni Rostock
  \end{columns}
\end{frame}

\begin{frame}{Kartografie}
  \begin{columns}
    \column{.6\textwidth}
	\includegraphics[width=\textwidth]{../images/landkarte.jpg}
	\column{.4\textwidth}
	Darstellung von Gewässern, Bewuchs und Schnee
  \end{columns}
\end{frame}

\begin{frame}{Maschinenbau}
  \begin{columns}
    \column{.6\textwidth}
	\includegraphics[width=0.9\textwidth]{../images/schnitt.jpg}
	\column{.4\textwidth}
	Vereinfachte technische Zeichnungen von Baugruppen\\
	Quelle: Gooch
  \end{columns}
\end{frame}

%\section{Wasserfarben}
\section{Wasserfarben}
\subsection{Einführung}
\begin{frame}{Einführung Wasserfarben}
  \begin{itemize}
    \item Moderne Wasserfarbentradition: zweite Hälfte des 18. Jahrhunderts
    \item Ermöglicht zahlreiche künstlerische Effekte
    \item Computer-Erzeugung durch zwei Ansätze (nach Curtis et al.): 
    realistische physikalische Simulation (hoher Aufwand) oder Annäherung durch 
    geeignete Filter (echtzeitgeeignet)
  \end{itemize}
\end{frame}

\begin{frame}{Materialien}
  \begin{itemize} 
    \item Papier: meist nicht aus Zellstoff, sondern Leinen bzw.\ Baumwolle;
    grobe Struktur, dadurch Wasserabsorption und Diffusion möglich
    \item Bestandteile der Wasserfarbe: in Wasser gelöstes Pigment, Bindemittel
    sowie Tensid (oberflächenaktiver Stoff)
    \item Pigment: festes Material, besteht aus einzelnen Partikeln; oftmals in
    Puderform; kann in das Papier einsickern und dort anhaften bzw.\ durch das
    Wasser transportiert werden
    \item Bindemittel: Pigment kann an das Papier anhaften
    \item Tensid: Wasser kann in das Papier eindringen
  \end{itemize}
\end{frame}

\begin{frame}{Haupt-Pinseltechniken}
  \begin{itemize}
    \item "`wet-on-wet painting"': Malen mit nassem, farbgetränkten Pinsel auf
    nassem Papier
    \item "`wet-on-dry painting"': Malen auf trockenem Papier
  \end{itemize}
\end{frame}

\begin{frame}{Effekte}
  \centering\pgfimage[width=1.0\textwidth]
    {../images/Curtis1997-real-watercolor-effects}
  
  \begin{itemize}
    \item \textsl{Dry-brush}: trockener Pinsel/Papier; unregelmäßige Lücken
    \item \textsl{Edge darkening}: Pigmente von innen n.\ außen; dunkle Kanten
    \item \textsl{Backruns}: Wasser läuft in nasse Region zurück
    \item \textsl{Granulation} u.\ \textsl{Separation}: körnige Struktur, Farbtrennung
    \item \textsl{Flow patterns}: freies Ausbreiten der Pinselstriche
    \item \textsl{Color glazing}: Übereinanderlegen von dünnen Farbschichten
  \end{itemize}
\end{frame}

\subsection{Computererzeugte Wasserfarben}
\begin{frame}{Simulation}
  \begin{itemize}
    \item Wichtig: nicht nur physikalische Eigenschaften, sondern auch
    künstlerische Effekte berücksichtigen!
    \item Bestandteile: Papierschicht, geordnete Menge von Farbschichten
    \item Für jede Farbschicht wird Flüssigkeitssimulation separat durchgeführt
    \item Zusätzlich: "`wet-area"'-Maske
    \item Nach der Simulation: Komposition mit Kubelka-Munk-Farbmodell erzeugt 
    finale Darstellung
  \end{itemize}
\end{frame}

\begin{frame}{Simulation: Papier}
  \begin{itemize}
    \item Vereinfachtes Papier-Modell: Textur wird als Höhenfeld und
    Flüssigkeitskapazitätsfeld abgebildet
    \item Höhe: $0 < h < 1$
    \item Steigung wird zur Berechnung der Fließgeschwindigkeit in der
    Simulation benutzt
    \item Kapazität: $c = h(c_{max} - c_{min}) + c_{min}$
  \end{itemize}
\end{frame}

\begin{frame}{Simulation: 3-Schichten-Modell}
  \begin{columns}
    \column{.6\textwidth}
	    \begin{enumerate}
	      \item "`shallow-water layer"': Wasser fließt über die Oberfläche und
	      transportiert ggf.\ Pigmente
	      \item "`pigment-deposition layer"': Pigmente werden hier abgelagert und
	      von hier weitertransportiert
	      \item "`capillary layer"': Absorbiertes Wasser diffundiert durch
	      Kapillarwirkung
	    \end{enumerate}

	\column{.4\textwidth}
		\pgfimage<3->[width=0.5\textwidth]{../images/Curtis1997-layer-shallow-water}\par
		\pgfimage<4->[width=0.5\textwidth]{../images/Curtis1997-layer-pigment-deposition}\par
		\pgfimage<5->[width=0.5\textwidth]{../images/Curtis1997-layer-capillary}
  \end{columns}
\end{frame}

\begin{frame}{Simulation: Haupt-Schleife}
  Eingabewerte: (Anfangswerte)
  \begin{itemize}
    \item Wet-Area-Maske $M$
    \item Geschwindigkeit des Wassers $u$ und $v$
    \item Wasserdruck $p$
    \item Pigmentkonzentration $g^k$
    \item Wassersättigung des Papiers $s$
  \end{itemize}
  
  \pause
  Simuliert werden über einen definierten Zeitraum: Wasser- und Pigmentbewegungen,
  Pigmentablagerungen, Kapillarfluss.
\end{frame}

\begin{frame}{Simulation: Rendering und Pigmente}
  \begin{itemize}
    \item Verwendung des Kubelka-Munk-Farbmodells
    \item Jedem RGB-Farbkanal jeden Pigments sind zwei Koeff. zugewiesen:
    \begin{itemize}
      \item Absorptionskoeff. $K$
      \item Streuungskoeff. $S$
    \end{itemize}
  \end{itemize}
  
  \begin{block}{Synthetische Pigmente}
    \pgfimage<5->[width=0.8\textwidth]{../images/Curtis1997-pigments}
  \end{block}
\end{frame}

\begin{frame}{Simulation: Vergleich der Ergebnisse}
  \centering\pgfimage[width=1.0\textwidth]
  {../images/Curtis1997-real-watercolor-effects}\par
  \centering\pgfimage[width=1.0\textwidth]
  {../images/Curtis1997-simulated-watercolor-effects}
\end{frame}

\begin{frame}{Anwendungen}
  \begin{itemize}
    \item Interaktives Malen: Benutzer kann Startzustand der Simulation
    "`malen"' durch Erstellung der Schichten mit Unterschichten für Pigmente,
    Wasser und die Wet-Area-Maske. Zusätzlich: globales Referenzbild sowie
    Papiertextur
    \item "`Watercolorization"': Konvertierung eines Farbbilds in eine
    Wasserfarben-Illustration.
    \item 3D-Szenen: Erweiterung der "`Watercolorization"'
    \item Aufbauende Ansätze: Darstellung in Echtzeit
  \end{itemize}
\end{frame}

\subsection{Echtzeit-Wasserfarben-Animationen}
\begin{frame}{Echtzeit-Animationen}
  \begin{itemize}
    \item Ansatz nach O. Deussen und T. Luft: keine exakte physikalische 
    Simulation, sondern Mittelweg zwischen Qualität und Rechenaufwand mit Ziel: 
    Echtzeit-Rendering von 3D-Szenen
    \item Aspekte: \textbf{Abstraktion und Vereinfachung},
    \textbf{Wasserfarben-Effekte}, \textbf{Licht und Schatten}
    \item Grobe Vorgehensweise: Erzeugung einzelner abstrakter 
    Wasserfarben-Layer, die mithilfe der Grafikkarte zusammengefügt werden
  \end{itemize}
\end{frame}

\subsubsection{Vorgehensweise}
\begin{frame}{Abstraktion und Vereinfachung}
  \begin{columns}[t]
    \column{0.6\textwidth}
      \begin{itemize}
        \item Jeder Layer enthält jeweils ein oder mehrere gleichartige Objekte
        \item Segmentieren der 3D-Szene anhand eindeutiger Identifikatoren
        (Ergebnis: "`intensity images"')
        \item Low-Pass-Filter erzeugt abstrakte, weiche Formen mit definierbarem
        Detaillierungsgrad
        \item Layer enthalten außerdem Farb- sowie Transparenzinformation
      \end{itemize}
    \column[T]{0.4\textwidth}
      \pgfimage<4->[width=1.0\textwidth]{../images/Luft2006-tree-segmentation}
  \end{columns}
\end{frame}

\begin{frame}{Formextraktion und Fließmuster}
  \begin{columns}[t]
    \column{0.6\textwidth}
      \begin{itemize}
        \item Form ergibt sich aus den Intensitäts-Werten der \textsl{intensity
        images}
        \item Fließmuster beschreibt das Pigment-Verhalten am Rand des Layers; 
        \textsl{wet-on-dry}: harte Kanten, \textsl{wet-on-wet}: weiche, 
        federartige Kanten
      \end{itemize}
    \column[T]{0.4\textwidth}
      \pgfimage<4->[width=1.0\textwidth]{../images/Luft2006-tree-shape-extraction-flow-pattern}
  \end{columns}
\end{frame}

\begin{frame}[label=current]{Edge darkening}
  \begin{columns}[t]
    \column{0.6\textwidth}
      \begin{itemize}
        \item Pigmente werden beim Trocknen zum Rand transportiert
        \item Imitation durch Anwendung eines Gauss-Filters, erzeugt weichen
        Intensitäts-Übergang an den Rändern
        \item Erzeugung eines Gradienten, der dem Wasserfarbenfluss folgt und
        die Transparenz an den Rändern verblassen lässt
      \end{itemize}
    \column[T]{0.4\textwidth}
      \pgfimage<4->[width=0.8\textwidth]{../images/Luft2006-tree-edge-darkening}
  \end{columns}
\end{frame}

\begin{frame}{Pigment-Granulation und Transparenz}
  \begin{columns}[t]
    \column{0.6\textwidth}
      \begin{itemize}
        \item Struktur des Papiers beeinflusst Wasserverlauf und erzeugt so Pigment-Granulation
        \item Imitation durch Verwendung einer zusätzlichen Papier-Textur
        \item Mit entsprechenden Texturen lassen sich auch ausgefranste Kanten erzeugen
      \end{itemize}
    \column[T]{0.4\textwidth}
      \pgfimage<3->[width=1.0\textwidth]{../images/Luft2006-tree-granulation-transparency}
  \end{columns}
\end{frame}

\subsubsection{Komposition von Wasserfarben-Layern}
\begin{frame}{Licht (1)}
  \begin{columns}
    \column{0.6\textwidth}
      \begin{itemize}
        \item Verwendung der Beleuchtungs-Informationen aus dem 3D-Modell zur
        Anpassung der erzeugten bzw.\ zur Erstellung neuer Layer
        \item Phong-Beleuchungsmodell, Komponenten: ambiente (hier konstant), 
        diffuse und spiegelnde Reflexion
        \item Diffuse und spiegelnde Reflexion werden auf zwei \textsl{lighting-maps} gerendert
      \end{itemize}
    \column{0.4\textwidth}
      \pgfimage<4->[width=1.0\textwidth]{../images/Luft2006-still-life-light-shadow}
  \end{columns}
\end{frame}

\begin{frame}{Licht (2)}
  \begin{columns}
    \column{0.6\textwidth}
      \begin{itemize}
        \item spiegelnde Reflexion: erzeugt hervorgehobene Region, wird erreicht
        durch Ausmaskieren der \textsl{intensity-images} vor Erstellung des
        Layers (Bild 3 und 4)
        \item diffuse Reflexion: Erzeugung zusätzlicher Layer, z.\,B.\
        Hervorhebung dunkler Stellen (Bild 3); weiterhin Anpassung aller Layer,
        zwei Farbwerte für unbeleuchtete und beleuchtete Regionen (Bild 2)
      \end{itemize}
    \column{0.4\textwidth}
      \pgfimage<2->[width=1.0\textwidth]{../images/Luft2006-still-life-light-shadow}
  \end{columns}
\end{frame}

\begin{frame}{Komposition}
  \begin{columns}[t]
    \column{0.6\textwidth}
      \begin{itemize}
        \item Layer werden mit Standard-Funktion zur Überblendung
        transparenter Objekte zusammengesetzt
        \item Farbe an der Position $(x,y)$ ergibt sich zu: $R_{rgb} = C_a \cdot
        C_{rgb} + (1-C_a) \cdot B_{rgb}$
        \item $C_{rgb}$ = Farbe, $C_a$ = Transparenz, $B_{rgb}$ = Hintergrundfarbe
      \end{itemize}
    \column[T]{0.4\textwidth}
      \pgfimage<4->[width=1.0\textwidth]{../images/Luft2006-tree-result}
  \end{columns}
\end{frame}

\begin{frame}{Demo}
  \begin{center}
	  \movie[autostart, height=6cm, width=8cm]{}{../tree_cover.avi}
  \end{center}
\end{frame}

% Schluss-Folie
\begin{frame}
  \frametitle{Danke.}
  \begin{columns}
    \column{.6\textwidth}
	\includegraphics[width=\textwidth]{../images/wir.jpg}
	\column{.4\textwidth}
	Fragen?
  \end{columns}
\end{frame}

\end{document}